\chapter{Etude Prealable}
\label{chap:etude-prealable}

\section{Introduction Generale}

Aujourd'hui, le monde evolue rapidement avec la transformation numerique. Les entreprises ne peuvent plus se limiter a utiliser des outils digitaux basiques. L'automatisation des processus, qui etait autrefois un avantage, est devenue une necessite. Mais desormais, cela ne suffit plus. Il est important de concevoir des systemes a la fois automatiques et intelligents, capables non seulement d'executer des taches repetitives, mais aussi d'analyser des donnees, de prendre des decisions pertinentes et d'ameliorer en continu les performances.

Le domaine des ressources humaines reflete bien cette evolution. Avec l'augmentation du nombre de candidatures et le besoin de rapidite et de precision, les methodes traditionnelles de recrutement montrent leurs limites. Le tri manuel des CV, l'evaluation des candidats, l'organisation des entretiens ou encore la communication avec les postulants prennent beaucoup de temps et ne sont pas toujours efficaces. Dans ce contexte, un simple systeme automatise ne suffit plus. Il devient necessaire de mettre en place une solution intelligente capable d'aider, voire d'ameliorer, la prise de decision humaine.

C'est dans ce cadre que s'inscrit ce projet de fin d'etudes. L'objectif est de concevoir et de developper une plateforme web innovante dediee a la gestion intelligente du recrutement. Cette solution permet d'importer et d'analyser automatiquement les CV, de les transformer selon des modeles predefinis et d'extraire les informations importantes. Elle propose egalement un module de gestion des offres d'emploi, qui peuvent etre creees manuellement ou generees a l'aide de l'intelligence artificielle.

En plus de ces fonctionnalites, le systeme integre un mecanisme de scoring automatique des candidats, base sur l'analyse de leurs competences et leur adequation avec les offres. Il permet aussi d'automatiser la communication avec les candidats, notamment grace a l'envoi d'emails personnalises, et facilite la planification des entretiens via une integration avec des outils comme Google Calendar. Ainsi, la solution ne se contente pas d'automatiser le recrutement, elle le rend plus simple, plus rapide et plus efficace.

Ce rapport presente les differentes etapes de realisation du projet. Il commence par une etude du contexte et des besoins, ainsi que les choix techniques adoptes. Ensuite, il decrit l'architecture du systeme et l'environnement de developpement. Les chapitres suivants detaillent les phases de realisation et les principales fonctionnalites mises en place. Enfin, le rapport se termine par une evaluation globale du systeme, incluant les tests effectues, les resultats obtenus et les ameliorations possibles a l'avenir.

\section{Presentation de l'organisme d'accueil}

\subsection{Presentation de Tritux}

Tritux est une entreprise specialisee dans le developpement de solutions numeriques et l'accompagnement des organisations dans leur transformation digitale. L'entreprise se concentre sur l'innovation technologique, l'automatisation des processus metier et l'integration de solutions intelligentes basees sur l'intelligence artificielle.

Tritux travaille principalement sur :
\begin{itemize}
    \item Le developpement d'applications web et mobiles,
    \item L'integration de systemes \acrfull{api} REST et microservices,
    \item L'automatisation de flux de travail via des outils comme n8n,
    \item L'exploitation de services \acrfull{ia} pour l'analyse et la structuration de donnees,
    \item L'infrastructure cloud et la containerisation Docker.
\end{itemize}

\subsection{Contexte du projet chez Tritux}

Dans le cadre de ses besoins internes et de ses processus de gestion des ressources humaines, Tritux souhaite disposer d'une plateforme capable de simplifier et d'optimiser le traitement des candidatures. L'objectif est de reduire le temps passe sur les taches manuelles et repetitives, tout en améliorant la qualite de l'analyse des profils de candidats.

Le projet est realise en tant que travail de fin d'etudes dans un environnement professionnel reel. Il s'inscrit dans une demarche agile avec des sprints bi-hebdomadaires, permettant une validation rapide des fonctionnalites et une adaptation continue aux besoins.

\section{Contexte et problematique}

\subsection{Situation actuelle du recrutement}

Actuellement, le processus de traitement des CV repose largement sur des operations manuelles :

\begin{itemize}
    \item \textbf{Reception des candidatures} : les CV arrivent par email, formulaires web ou portails de recrutement, sans centralisation,
    
    \item \textbf{Tri et classement manuel} : chaque candidature doit etre triee et classee selon des criteres heterogenes,
    
    \item \textbf{Lecture et extraction d'informations} : extraction manuelle des donnees pertinentes (formation, experiences, competences),
    
    \item \textbf{Reformatage} : adaptation du CV aux standards de l'entreprise (police, mise en page, sections),
    
    \item \textbf{Communication separee} : envoi d'emails individuels et peu personnalises aux candidats,
    
    \item \textbf{Suivi fragmentaire} : absence de suivi centralise de l'etat de chaque candidature.
\end{itemize}

\subsection{Problemes identifies}

Cette approche manuelle presente plusieurs limitations critiques :

\begin{itemize}
    \item \textbf{Consommation de ressources} : les equipes RH consacrent beaucoup de temps a des taches repetitives et peu valorisantes,
    
    \item \textbf{Risque d'erreurs} : les donnees extraites manuellement peuvent etre incompletes, incoherentes ou erronees,
    
    \item \textbf{Manque de standardisation} : les CV conservent des formats heterogenes rendant l'analyse comparative difficile,
    
    \item \textbf{Scalabilite limitee} : le traitement manuel ne peut pas suivre une augmentation du volume de candidatures,
    
    \item \textbf{Absence d'intelligence d'entreprise} : pas d'analyse crossfunctionnelle des profils et des competences,
    
    \item \textbf{Communication inefficace} : communications non-personnalisees et demarches ad-hoc de notification.
\end{itemize}

\section{Objectifs du projet}

\subsection{Objectif general}

Concevoir et developper une plateforme web intelligente et modulaire de gestion du recrutement permettant l'import, l'analyse automatique, la transformation et la gestion des candidatures avec orchestration intelligente des flux de travail.

\subsection{Objectifs specifiques}

Les objectifs specifiques sont :

\begin{enumerate}
    \item \textbf{Centralisation des candidatures} : creer une plateforme web centralisee permettant l'import de CV depuis differentes sources,
    
    \item \textbf{Extraction intelligente} : extraire automatiquement les informations pertinentes des CV (nom, competences, experiences, formation) avec precision,
    
    \item \textbf{Structuration des donnees} : normaliser et structurer les donnees extraites dans un format standardise exploitable,
    
    \item \textbf{Transformation parametree} : permettre la transformation des CV selon differents templates de presentation,
    
    \item \textbf{Gestion des offres d'emploi} : gerer et diffuser les offres d'emploi avec support multi-langue,
    
    \item \textbf{Automatisation de la communication} : automatiser la notification des candidats avec messages personnalises,
    
    \item \textbf{Orchestration des workflows} : integrer l'orchestration de workflows pour coordonner les differentes etapes du processus,
    
    \item \textbf{Fournir une \acrfull{ui}} : offrir une interface utilisateur ergonomique et intuitive pour les gestionnaires RH.
\end{enumerate}

\section{Etude de l'existant}

\subsection{Processus actuel}

Le processus de gestion des CV dans un contexte traditionnel suit generalement le schema suivant :

\begin{enumerate}
    \item Reception des CV par email ou formulaire,
    \item Tri manuel selon differents criteres (experience, competences, localisation),
    \item Lecture individuelle et synthèse manuelle,
    \item Adaptation du CV a un format standardise,
    \item Communication individuelle avec les candidats selectionnes,
    \item Archivage dans un dossier partagé ou un systeme peu structuré.
\end{enumerate}

\subsection{Limites de l'existant}

Le processus actuel presente plusieurs limitations majeures :

\begin{itemize}
    \item \textbf{Manque de centralisation} : les CV et les donnees sont dispersés et peu accessibles de manière unifiée,
    
    \item \textbf{Absence de standardisation des donnees} : chaque CV conserve sa propre structure rendant la comparaison difficile,
    
    \item \textbf{Manque d'automatisation} : les taches repetitives ne sont pas automatisees, generant des couts en ressources humaines,
    
    \item \textbf{Absence d'intelligence d'entreprise} : aucune analyse crossfunctionnelle ou scoring automatique des candidats,
    
    \item \textbf{Suivi insuffisant} : le statut des candidatures et les historiques d'interactions ne sont pas centralises,
    
    \item \textbf{Absence d'orchestration} : les differentes etapes ne sont pas coordonnees et de nombreuses operations manuelles persistent.
\end{itemize}

\section{Etude technique de l'existant}

Aucune infrastructure technologique n'etait en place pour supporter le processus de gestion des candidatures. Par consequent, le projet a necessite la creation d'une plateforme completement nouvelle integrant:

\begin{itemize}
    \item Un backend de gestion et d'orchestration des donnees,
    \item Une interface frontend ergonomique,
    \item Des services de traitement intelligents,
    \item Une orchestration de workflows,
    \item Une base de donnees centralisee.
\end{itemize}

\section{Solution proposee}

\subsection{Vue generale de la solution}

La solution proposee est une plateforme web modulaire et orchestree composee de :

\begin{itemize}
    \item \textbf{Frontend React} : interface utilisateur moderne et reactive permettant l'import, la visualisation et la gestion des candidatures,
    
    \item \textbf{Backend Spring Boot} : \acrfull{api} REST assurant la persistance des donnees, l'orchestration des operations et la gestion des entites metier,
    
    \item \textbf{Service Python FastAPI} : service specialise dans l'extraction du texte des documents PDF et DOCX avec la bibliotheque Docling,
    
    \item \textbf{Orchestrateur n8n} : coordination intelligente des workflows et integration avec les services externes (Gemini API),
    
    \item \textbf{Base de donnees PostgreSQL} : stockage relationnels des donnees de candidats, offres d'emploi et templates,
    
    \item \textbf{Conteneurisation Docker} : deployement et scalabilite de l'infrastructure.
\end{itemize}

\subsection{Flux de traitement}

Le flux general de traitement suit les etapes suivantes :

\begin{enumerate}
    \item \textbf{Import et stockage} : l'utilisateur importe un CV (DOCX ou PDF) via l'interface React. Le fichier est stocke sur le serveur,
    
    \item \textbf{Extraction du texte} : le service Python utilise Docling pour extraire le texte brut du document,
    
    \item \textbf{Orchestration n8n} : n8n reçoit le texte et construit un prompt structuré pour l'API Gemini,
    
    \item \textbf{Analyse par Gemini} : Gemini analyse le texte et retourne un objet JSON structuré contenant nom, experiences, competences, formations, langues, etc.,
    
    \item \textbf{Validation Pydantic} : le service Python valide et normalise la reponse grace au modele Pydantic \texttt{CandidateProfile},
    
    \item \textbf{Stockage en base} : les donnees normalisees sont stockees dans PostgreSQL (table \texttt{candidates}) avec le JSON structuré,
    
    \item \textbf{Affichage frontend} : l'interface React affiche le profil du candidat et permet la visualisation des donnees extraites,
    
    \item \textbf{Transformation templates} : l'utilisateur peut appliquer un template de presentation pour generer un CV reformate (DOCX).
\end{enumerate}

\subsection{Entites metier}

La solution manipule trois entites metier principales :

\begin{enumerate}
    \item \textbf{Candidate} : represente un candidat avec ses donnees extraites du CV, stockee en table PostgreSQL \texttt{candidates},
    
    \item \textbf{JobOffer} : represente une offre d'emploi avec support multi-langue (FR/EN), stockee en table \texttt{job\_offers},
    
    \item \textbf{Template} : represente un modele de presentation de CV utilisable pour transformer un profil candidat.
\end{enumerate}

\subsection{Modules de l'interface}

L'interface React propose les modules fonctionnels suivants :

\begin{itemize}
    \item \textbf{Dashboard} : vue d'ensemble du processus de recrutement,
    \item \textbf{Profiles} : gestion et visualisation des candidats,
    \item \textbf{Templates} : gestion des modeles de presentation,
    \item \textbf{Transformation} : application d'un template a un profil pour generer le CV transforme,
    \item \textbf{JobOffers} : gestion des offres d'emploi,
    \item \textbf{Chat} : interaction avec un assistant IA pour des operations intelligentes.
\end{itemize}

\section{Environnement technologique adopte}

\subsection{Stack technologique}

Le projet utilise une stack technologique moderne et eprouvee :

\begin{itemize}
    \item \textbf{Frontend} : React 18.3, React Router 6.23, Tailwind CSS 3.4, Axios 1.7,
    
    \item \textbf{Backend} : Spring Boot 3.2.5, Java 21, Spring Data JPA, Spring Web,
    
    \item \textbf{Base de donnees} : PostgreSQL 14+,
    
    \item \textbf{Service Python} : FastAPI 0.111, Docling 2.5, Pydantic 2.7, python-docx 1.1,
    
    \item \textbf{Orchestration} : n8n 0.x+,
    
    \item \textbf{Services IA} : Google Gemini \acrfull{api},
    
    \item \textbf{Containerisation} : Docker et Docker Compose.
\end{itemize}

\subsection{Ports et architecture deploiement}

L'architecture deploiement distribue utilise les ports suivants :

\begin{table}[!ht]
\centering
\begin{tabular}{p{5cm}p{2cm}p{4cm}}
\toprule
\textbf{Service} & \textbf{Port} & \textbf{Description} \\
\midrule
Frontend React & 3000 & Interface utilisateur web \\
Backend Spring Boot & 9091 & API REST et orchestration \\
Service Python & 8001 & Extraction et traitement de documents \\
n8n Orchestrateur & 5678 & Workflows et orchestration intelligente \\
PostgreSQL & 5432 & Base de donnees relationnelles \\
\bottomrule
\end{tabular}
\caption{Ports et services de la plateforme}
\label{tab:ports-services}
\end{table}

\section{Methodologie de travail}

\subsection{Approche Scrum Agile}

Le projet est realise selon la methodologie Scrum Agile, qui favorise l'iterativite, la collaboration et la livraison incrementale de fonctionnalites. Cette approche est particulierement adaptee aux projets de developpement logiciel modernes necessitant une adaptation rapide aux changements.

Les principes fondamentaux de Scrum adoptes sont :

\begin{itemize}
    \item \textbf{Iterativite} : le travail est organise en sprints de courte duree (2 semaines) permettant des cycles rapides de developpement et de feedback,
    
    \item \textbf{Transparence} : communication reguliere entre l'equipe de developpement et les parties prenantes via les reunions de sprint,
    
    \item \textbf{Inspection et adaptation} : evaluation reguliere des progres et ajustement des priorites selon les retours,
    
    \item \textbf{Livraison incrementale} : chaque sprint livre une portion de fonctionnalites operationnelles et testees.
\end{itemize}

\subsection{Organisation des sprints}

Le projet a ete organise en 8 sprints bi-hebdomadaires (16 semaines) suivant cette structure :

\begin{enumerate}
    \item \textbf{Sprint 1-2} : setup de l'environnement, architecture et design initial des bases de donnees,
    
    \item \textbf{Sprint 3-4} : developpement du module d'import des CV et integration du service Python,
    
    \item \textbf{Sprint 5-6} : developpement de l'extraction de donnees, orchestration n8n et validation Pydantic,
    
    \item \textbf{Sprint 7-8} : developpement de la gestion des offres d'emploi, templates et interface de transformation,
    
    \item \textbf{Sprint 9-10} : finalisation des modules, tests et optimisations,
    
    \item \textbf{Sprint 11-12} : integration des retours utilisateurs et ameliorations ergonomie,
    
    \item \textbf{Sprint 13-14} : tests de performance et scalabilite,
    
    \item \textbf{Sprint 15-16} : préparation du deploiement et documentation finale.
\end{enumerate}

\subsection{Ceremonies Scrum}

Les ceremonies Scrum suivantes ont ete etablies :

\begin{itemize}
    \item \textbf{Sprint Planning} : reunion de 2-4 heures au debut de chaque sprint pour selectionner les elements du backlog et estimer le travail,
    
    \item \textbf{Daily Standup} : reunion courte (15 min) chaque jour pour synchroniser l'equipe sur les taches et identifier les obstacles,
    
    \item \textbf{Sprint Review} : demonstration des fonctionnalites realisees a la fin de chaque sprint aux parties prenantes et clients,
    
    \item \textbf{Sprint Retrospective} : reflexion d'equipe sur les ameliorations du processus et de la collaboration.
\end{itemize}

\subsection{Roles Scrum}

L'equipe Scrum est composee de :

\begin{itemize}
    \item \textbf{Product Owner} : represente les interets des utilisateurs finaux et prioritaires, gère le backlog produit,
    
    \item \textbf{Scrum Master} : facilite le processus Scrum, elimine les obstacles techniques ou organisationnels,
    
    \item \textbf{Equipe de developpement} : composee de developpeurs (backend et frontend), architecte et testeurs responsables de la realisation.
\end{itemize}

\section{Conclusion}

Ce chapitre a presente le contexte de realisation du projet, les limites de l'existant, les objectifs cibles et la vision globale de la solution. 

La plateforme Tritux RH represente une reponse intelligente aux defis du recrutement modern, integrant une architecture modulaire multi-service (React frontend, Spring Boot backend, FastAPI microservice, orchestration n8n) pour automatiser et optimiser l'ensemble du processus de gestion des candidatures.

Le chapitre suivant detaille les besoins fonctionnels et non fonctionnels qui ont guide la conception de cette solution.